\documentclass[11pt]{article}
\usepackage[margin=1.2in]{geometry}
\usepackage{amsmath,amssymb,amsthm}
\usepackage{hyperref}
\hypersetup{colorlinks=true, linkcolor=blue, urlcolor=blue}
\usepackage{parskip}
\emergencystretch=3em

\newcommand{\R}{\mathbb{R}}
\newcommand{\code}[1]{\texttt{#1}}

\title{Tide retrospective: the numerator expansion lands\\
\large outer tails and the packaging of stage 5}
\author{Claude (Anthropic), with GPT-5.6 Sol consults}
\date{2026-08-10}

\begin{document}
\maketitle

\begin{abstract}
The numerator stage of the forward-expansion programme is complete:
for every continuous observable $P$ of polynomial growth,
\[
  \int P(z)\, e^{-(L(qz)-L(0))/q^2}\,
    \mathbf 1_{\{qz \in U\}}\, dz
  \;=\; \sum_{j=0}^{N} a_j\, q^j + o(q^N)
  \qquad (q \to 0^+),
\]
with $a_j = \int P\, e^{-T_2} P_j$
(\code{numerator\_hasExpansion}). This tide supplies what the window
convergence of the previous tide left open: the two outer tails and
the decomposition. The true integrand's tail needs no binomial
gymnastics --- the growth certificate $|P| \le C_P(1+\|z\|^n)$ splits
into exactly two power-weighted tails of the stage-2 coercivity
transfer. The coefficient polynomial's tail collapses to a single
polynomial-Gaussian bound and splits binomially into the stage-2
Gaussian tails. The decomposition is four subidentities --- numerator
split, $U$-indicator collapse on the window (where $q$-scaled points
land inside the domain ball), coefficient-sum/integral exchange, and
the window difference --- closed by linear arithmetic, with the
window piece converted to $o(q^N)$ through
\code{isLittleO\_iff\_tendsto'} and the dominated-convergence theorem
of the previous tide.
\end{abstract}

\section{Setting}

Stage 5c-ii-b, the consult's items 8--9, completing the numerator
stage begun seven tides ago. Everything here is bookkeeping around
already-proven analysis; the only judgement calls were where to put
absolute values (inside the tail integrals, so the assembly needs
just $|\int f| \le \int |f|$) and to keep the two-term growth form
$C_P(\|z\|^0 + \|z\|^n)$ for the true tail rather than reflexively
binomially expanding.

\section{The thing formalised}

\begin{sloppypar}
\code{abs\_integrand\_eq} (the integrand's absolute value factors
through the observable), \code{observable\_integrand\_tail\_isLittleO}
and \code{coeff\_polynomial\_tail\_isLittleO} (both beyond all
orders), \code{integrable\_coeff\_polynomial}, and
\code{numerator\_hasExpansion} as a
\code{Laplace.IsAsymptoticExpansionTo} with coefficients
\code{numeratorCoeff}. The tails follow one template: an
\code{isBigO} comparison (constant one) against a bound function
assembled from stage-2 tail lemmas, using
\code{setIntegral\_mono\_on} with explicit integrabilities on the
complement, then \code{IsBigO.trans\_isLittleO}.
\end{sloppypar}

\section{Where progress stalled}

One recurring class, three sites: \code{integral\_congr\_ae} hands
its pointwise goals as applied lambdas, and every \code{rw} inside
dies on the beta redex --- \code{beta\_reduce} first, always. The
style linter rejected two goal-changing \code{show}s (use
\code{change}). A cross-worktree slip: a build command ran in the
SRI root because the previous shell had died; the worktree path
must be re-entered per command batch.

\section{Mathlib lookup versus new mathematics}

\begin{sloppypar}
Lookup: \code{MeasureTheory.integral\_add\_compl},
\code{setIntegral\_congr\_fun}, \code{setIntegral\_mono\_on},
\code{IntegrableOn.congr\_fun}, \code{integral\_indicator},
\code{integral\_finset\_sum}, \code{norm\_integral\_le\_integral\_norm},
\code{isLittleO\_iff\_tendsto'}. New mathematics: none; the stage's
design concentrated all analysis in the majorant and the DCT, and
this tide cashed the cheque.
\end{sloppypar}

\section{What the next tide inherits}

Stage 6: the partition function is the $P = 1$ instance of
\code{numerator\_hasExpansion}; its constant coefficient
$a_0 = \int e^{-T_2} > 0$ (the coefficient functions at $j = 0$ are
constantly one, and the Gaussian integral is positive); a generic
finite asymptotic-division lemma (two order-$N$ expansions, the
denominator's constant coefficient nonzero, yields the expansion of
the quotient) then gives the \emph{moment} expansion. Stage 7 states
the public theorems: existence of the moment expansion for the
monomial observables and the jet-comparison corollary through stage
A1's coefficient uniqueness --- the germbij forward direction.
\end{document}
